% !TeX root = main.tex
\chapter*{ПЕРЕЛІК УМОВНИХ ПОЗНАЧЕНЬ, СИМВОЛІВ, ОДИНИЦЬ, СКОРОЧЕНЬ І ТЕРМІНІВ}
\phantomsection
\addcontentsline{toc}{chapter}{ПЕРЕЛІК УМОВНИХ ПОЗНАЧЕНЬ, СИМВОЛІВ, ОДИНИЦЬ, СКОРОЧЕНЬ І ТЕРМІНІВ}

\begingroup
\small
% Longtable створює посилальну мітку навіть без назви. Службове значення
% лічильника зберігає унікальну мітку table.0.1 для першої нумерованої таблиці.
\setcounter{table}{-1}
\begin{longtable}{>{\bfseries}p{0.22\textwidth}p{0.70\textwidth}}
\multicolumn{2}{l}{\textbf{Скорочення та назви форматів}}\\
AUC & площа під кривою залежності повноти від частки хибнопозитивних рішень;\\
Acc & частка правильних рішень;\\
DCT & дискретне косинусне перетворення;\\
EXIF & формат службових даних зображення;\\
$F_1$ & гармонійне середнє точності позитивного виявлення та повноти;\\
FN & кількість хибнонегативних рішень;\\
FNR & частка хибнонегативних рішень;\\
FP & кількість хибнопозитивних рішень;\\
FPR & частка хибнопозитивних рішень;\\
HEVC & стандарт кодування відеоданих;\\
IEND & завершальний блок структури PNG;\\
IPTC & стандарт полів службових даних мультимедійних об’єктів;\\
JPEG & формат стиснення та зберігання зображень;\\
LSB & найменш значущий біт;\\
MP4 & мультимедійний контейнер на основі ISO Base Media File Format;\\
PNG & формат растрових зображень без втрат;\\
PR & крива залежності точності позитивного виявлення від повноти;\\
Prec & точність позитивного виявлення;\\
PSNR & пікове відношення сигналу до шуму;\\
Rec & повнота виявлення;\\
RIFF & контейнерний формат обміну блоковими даними;\\
ROC & крива залежності повноти від частки хибнопозитивних рішень;\\
SHA-256 & криптографічна геш-функція з 256-бітовим результатом;\\
TN & кількість правильно визначених контрольних об’єктів;\\
TP & кількість правильно визначених контрольовано змінених об’єктів;\\
ViT & модель перетворювача для аналізу зображень;\\
WebM & відкритий мультимедійний контейнер на основі EBML;\\
WebP & формат растрових зображень зі стисненням із втратами або без втрат;\\
XMP & розширювана платформа службових даних.\\
\addlinespace
\multicolumn{2}{l}{\textbf{Основні математичні символи}}\\
$X_{mm}$ & початковий мультимедійний об’єкт як незмінна послідовність байтів;\\
$X_c$ & сукупність статичних представлень об’єкта;\\
$X_A$ & повний вхід моделі з даними, маскою доступності та супровідними відомостями;\\
$X_b$ & послідовність причинно пов’язаних подій;\\
$Z_s$, $Z_v$, $Z_b$ & представлення відповідно структурних ознак, спектральних і піксельних ознак та ознак подій і поведінки;\\
$Z_c$ & маскована числова проєкція доступних статичних ознак;\\
$H$ & статичний опис об’єкта з оцінками, координатами, маскою та відомостями про походження;\\
$Z$, $Z^{\partial}$ & повне та неповне інтегровані представлення;\\
$m_c$, $m_A$, $m_{bq}$ & маски доступності статичних складових, трьох груп даних і спільної поведінкової оцінки відповідно;\\
$L_s$, $L_v$ & байтові та просторово-часові координати ознак;\\
$\rho_s$, $\rho_v$, $\rho_b$, $\rho_A$ & часткові та інтегральна оцінки ризику;\\
$Q_A$ & оцінка причинно-фазової відповідності поведінкових ознак;\\
$d(e)$ & складена ознака допустимості події до поведінкового узгодження;\\
$C_M$ & верхня оцінка часової складності методу поведінкового узгодження;\\
$N_J$, $T$, $n_r$ & кількість записів початкового журналу, допущених подій і правил, що перевіряються для однієї події;\\
$F_C$, $F_I$, $F_D$ & функції маскованого статичного об’єднання, інтеграції та прийняття рішення;\\
$\mathbf 1[\cdot]$ & характеристична функція виконання умови;\\
$\mathrm R$ & множина дійсних чисел;\\
$\mathrm I_A$, $n_m$ & множина доступних груп ознак і кількість її елементів;\\
$\mathrm K$, $\mathrm K_j$ & множина класів аналітичного стану та окремий клас;\\
$\mathrm T_D$ & множина параметрів прийняття рішення;\\
$\mathrm P^{n_A-1}$ & множина нормованих векторів класових оцінок;\\
$\mathrm E$ & множина структурованих пояснень;\\
$\mathrm R_{\mathrm{app}}$ & множина правил, застосовних до наявних подій;\\
$\theta$, $\Theta$ & параметри окремого перетворення та складена версована конфігурація відповідно;\\
$S_A$ & аналітичний стан об’єкта без технологічної дії;\\
$K_I$ & вторинний розвідувальний рейтинг конфігурації;\\
$\bot$ & недоступне або невизначене значення, яке не є числовим нулем;\\
$\varnothing$ & порожня, але доступна для спостереження множина або послідовність.\\
\addlinespace
\multicolumn{2}{l}{\textbf{Одиниці вимірювання}}\\
дБ & децибел;\\
мс & мілісекунда;\\
МіБ & мебібайт, $2^{20}$ байтів;\\
об’єкт/с & кількість опрацьованих об’єктів за секунду.\\
\end{longtable}
\endgroup
% The unnumbered list above must not shift numbering of the first content table.
\setcounter{table}{0}

\clearpage
